% New English translation from checked cote 25, site 88 (PDF page 89).
% AI-translated by GPT-6 Astra (OpenAI, ChatGPT Pro mode).
\scanpage{25/88}
\mspage{\hand{82}}

Grass$^{m}$(P), the dimension of $\underline{X}^{(m)}_{\hand{\xi}}$ is at least dim X-m; this holds
% NOTE: the right edge of the scan cuts the ends of the first two source lines.
in particular when $\hand{\xi}$ comes from a point of Q(C$_{\hand{\unsure{\eta|y}}}$). For the inequality
in the opposite direction, note that (working over the field k' = k($\hand{\xi}$)),
since $C_{\hand{\unsure{\eta|y}}\,k'} \hand{\subset} L_{\hand{\xi}}$ is \var{\xout{of}}{}{typewriter strike-out}{ega25_25_088-v0001} a hyperplane of $L_{\hand{\xi}}$, if the dimension of
$\underline{X}^{(m)}_{\hand{\xi}} = Xx_{P}L_{\hand{\xi}}$ were \hand{$\geqslant$} dim X - m + 1, then that of $\underline{X}^{(m+1)}_{\hand{\unsure{\eta|y}}}$ =
$Xx_{P}C_{\hand{\unsure{\eta|y}}}$ would be \hand{$\geqslant$} dim X - m \var{\xout{+}}{}{typewriter strike-out}{ega25_25_088-v0002}, (since the base change k($\hand{\unsure{\eta|y}}$)$\hand{\to}$k'
transforms the latter prescheme into $(Xx_{P}L_{\hand{\xi}})x_{L_{\hand{\xi}}}(C_{\hand{\unsure{\eta|y}}\,k'})$). But on the
contrary, we have \var{\xout{dim $Xx_PC$}}{}{typewriter strike-out}{ega25_25_088-v0004} dim $\underline{X}^{(m+1)}$ = dim X - m -1 \var{\xout{(with the convention that}}{}{typewriter strike-out}{ega25_25_088-v0005}
\var{\deleted{(at least assuming dim f(X) \hand{$>$} m)}}{}{deletion}{ega25_25_088-v0006}, by no. 2 (revised in no. 10).
The \var{\inserted{case b)}}{case b)}{insertion}{ega25_25_088-v0007} \var{\deleted{case where dim f(X)=m can also be dealt with by case b), which}}{}{deletion}{ega25_25_088-v0008} is treated
similarly: if we had dim $Xx_{P}L$ \hand{$\geqslant$} d\hand{+1}, \var{\xout{then the image of}}{}{typewriter strike-out}{ega25_25_088-v0009} or
equivalently \var{\xout{\unsure{$f(X)\ f_{k(\ )}$}}}{}{typewriter strike-out}{ega25_25_088-v0010} $f_{k'}(X_{k'}) \hand{\cap} L_{\hand{\xi}}$ of dimension \hand{$\geqslant$}\var{\over{0}{1}}{1}{writing over}{ega25_25_088-v0011},
then by the same argument as \var{\xout{above}}{}{typewriter strike-out}{ega25_25_088-v0012} above we would have
$\underline{X}^{(m+1)}_{\hand{\unsure{\eta|y}}} \hand{\neq} \hand{\emptyset}$, contrary to what we pointed out before 14.1.

\uline{Corollary} 14.2. Suppose X of dimension m and f:X$\hand{\to}$P finite\var{\inserted[typed]{(resp quasi-finite)}}{ (resp quasi-finite)}{insertion}{ega25_25_088-v0013}; then
the morphism $p_{C_{\hand{\unsure{\eta|y}}}} : X_{k(\hand{\unsure{\eta|y}})} \hand{\unsure{\longrightarrow|\relbar}} Q(C_{\hand{\unsure{\eta|y}}})$ is finite surjective (resp. quasi-finite
dominant).

Indeed, this morphism is \var{\inserted[typed]{quasi-}}{quasi-}{insertion}{ega25_25_088-v0014}finite, and since dim $X_{\hand{k(\unsure{\eta|y})}}$ = dim Q (C$_{\hand{\unsure{\eta|y}}}$), it is
dominant; \var{\xout{finally hence surjective}}{}{typewriter strike-out}{ega25_25_088-v0015} if f is finite, $p_{C_{\hand{\unsure{\eta|y}}}}$ is moreover finite,
\var{\xout{\unsure{hence cl}}}{}{typewriter strike-out}{ega25_25_088-v0016} because it is proper, hence surjective since dominant.
\uline{Corollary} 14.3. Under the conditions of 14.1 a), if X is Cohen–Macaulay,
then the morphism $p_{C} : \overset{\hand{\sim}}{X}(C_{\hand{\unsure{\eta|y}}}) \hand{\unsure{\longrightarrow|\relbar}} Q(C_{\hand{\unsure{\eta|y}}})$ is a Cohen–Macaulay
morphism, and a fortiori \uline{flat}.

For the proof, cf the remark at the top of page 21, before 5.6.\var{\inserted{which gives a stronger result (to be included in 14.3), namely that the $\varphi^{(m)}_{\xi}$ for $\xi\in Q(C_{\unsure{\eta|y}})$ are \unsure{$F_{k(\xi)}$|$F_{\xi}$}-regular.}}{ which gives a stronger result (to be included in 14.3), namely that the $\varphi^{(m)}_{\xi}$ for $\xi\in Q(C_{\unsure{\eta|y}})$ are \unsure{$F_{k(\xi)}$|$F_{\xi}$}-regular.}{insertion}{ega25_25_088-v0017}
This \var{\over{c}{C}}{C}{writing over}{ega25_25_088-v0018}orollary should be modulated, but for simplicity we may
assume f quasi-finite: if F is a Cohen–Macaulay Module on X, and if
for \uline{every} irreducible component Z of Supp F we have dim Z \hand{$\geqslant$} m, then
$\overset{\hand{\sim}}{F}(C_{\hand{\unsure{\eta|y}}})$ is CM (and a fortiori flat) relative to Q(C$_{\hand{\unsure{\eta|y}}}$).

Note that we cannot replace,\var{\xout{\unsure{the hy}}}{}{typewriter strike-out}{ega25_25_088-v0019} to obtain the
\tnote{25-088-1}{The source's unindexed expression underlined X with superscript (m+1), and the references n°2 and n°10, are retained. Modulated renders modulé, followed in the source by the Module formulation. The underlined u within toute is represented by underlining the corresponding English word every. The canceled stray g in transforme gives no English correction.}

\var{\margin{88}}{\placedmargin{88}}{marginal note placed}{ega25_25_088-v0020}
